\scanpage{186}
\mspage{II.34}
% archive sheet number: 185
% block: top, lines 1-2 | text
forment un système projectif de groupes algébriques, noté $\widetilde G=\pi^{\mathrm{alg}}=(\widetilde G_E)$. On définit de façon triviale

% block: upper third, lines 3-6 | diagram
\diagram{\begin{tikzcd}[ampersand replacement=\&,column sep=large]
\operatorname{Mod}(\pi^{\mathrm{alg}})\arrow[d,no head] \& \operatorname{Mod}_{\mathbb Q_\ell}(\pi)\arrow[l,maps to,"\simeq"']\\
\begin{gathered}\text{ensemble de}\\\text{représentations}\\\text{\uline{algébriques}}\\\text{de }\pi^{\mathrm{alg}}\end{gathered}\&
\end{tikzcd}}

% block: upper to middle third, lines 7-8 | text
qui est une équivalence de catégories. On aura d'ailleurs

% block: middle third, line 9 | display
\[
\pi^{\mathrm{alg}}/(\pi^{\mathrm{alg}})^0\simeq\pi
\]

% block: middle third, line 10 | text
ce qui signifie que, si $\pi^{\mathrm{alg}}=(\underset{\text{groupe alg}}{\widetilde G_\alpha})$, alors

% block: middle third, line 11 | display
\[
\pi\unsure{\simeq|\approx}(\widetilde G_\alpha/\widetilde G_\alpha^0)
\]

% block: middle third, lines 12-13 | text
(isom \unsure{canoniques|canonique}), où bien sûr on identifie les groupes finis ordinaires à des groupes alg. sur $\mathbb Q_\ell$.
\margin{\unsure{C algé.|C algri.|G alg.}}

% block: middle to lower third, lines 14-15 | text
\deleted{6.6} Au système projectif $\pi^{\mathrm{alg}}=(\widetilde G_\alpha)$ est associé une pro-Algèbre de Lie

% block: lower third, line 16 | display
\[
\deleted{\unsure{\mathfrak g}}\widetilde{\mathfrak g}=\operatorname{Lie}^{\mathrm{alg}}(\pi)^{\deleted{\unsure{\mathrm{an}}}}=\operatorname{Lie}(\widetilde G^{\deleted{\unsure{\mathrm{alg}}}})=(\operatorname{Lie}(\widetilde G_\alpha)),
\]

% block: lower third, line 17 | text
et on \deleted{\unsure{a}} a un foncteur évident

% block: lower third, lines 18-20 | diagram
\diagram{\begin{tikzcd}[ampersand replacement=\&,column sep=large]
\operatorname{Mod}_{\mathbb Q_\ell}(\pi)\arrow[r,phantom,"\simeq"] \& \operatorname{Mod}(\deleted{\pi^{\mathrm{alg}}}\widetilde G)\arrow[r] \& \operatorname{Mod}(\widetilde{\mathfrak g})\deleted{\unsure{\operatorname{Mod}(\mathfrak g)}}\arrow[d,no head]\\
\&\&\begin{gathered}\varinjlim\text{ des cat. de}\\\text{modules sur les }\widetilde{\mathfrak g}_\alpha\\\text{de dim fin. sur }\mathbb Q_\ell\end{gathered}
\end{tikzcd}}

% block: lower third, line 21 | text
\deleted{\unsure{on obtient encore}}\inserted{et}

% block: bottom, lines 22-26 | diagram
\diagram{\begin{tikzcd}[ampersand replacement=\&,column sep=large]
\varinjlim_{\unsure{U_i|H_i}}\operatorname{Mod}_{\mathbb Q_\ell}(U_i)\arrow[r,"\simeq"]\arrow[d,no head]
\&\operatorname{Mod}(\widetilde G^0)\arrow[r]\arrow[d,no head]\&\operatorname{Mod}(\widetilde{\mathfrak g})\\
\begin{gathered}\text{parcourt les}\\\text{s-groupes}\\\text{ouverts de }\unsure{\pi|U_i}\end{gathered}
\&\begin{gathered}\text{syst. projectif}\\\text{des composantes}\\\text{connexes }G_\alpha^0\\\text{des }G_\alpha\end{gathered}\&
\end{tikzcd}}
